\item The sequence below (where $\tau_1(a)=a\widehat\otimes1$ and $\pi$ is the projection $A\to A/J_B$) is exact:
\[
\xymatrix@C=45pt{B\ar[r]&A\ar@<.6ex>[r]^{\tau_1}\ar@<-.6ex>[r]_{(\pi\widehat\otimes\id_A)\Delta_A}&A\widehat\otimes_k(A/J_B),}\tag{*}
\]
\ie $B$ is the set of all $x\in A$ such that $\Delta_A(x)-x\widehat\otimes1$ belongs to $A\widehat\otimes_k J_B$.
\end{enumeratei}
In this case, $H=\Spf(A/J_B)$ is a formal subgroup of $G$, and the sequence below (where $\lambda$ is the restriction to $G\times H$ of the multiplication of $G$) is exact:
\[
\xymatrix@C=38pt{G\times H\ar@<.6ex>[r]^{\operatorname{pr}_1}\ar@<-.6ex>[r]_\lambda&G\ar[r]&\Spf(B),}\tag{**}
\]
\ie $\Spf(B)$ is isomorphic to $G/H$.
\end{proposition}
% typo?: pi tensor id_A in (*) is retained, although its target has A/J_B on the right.
Put $\overline A=A/J_B$ and $H=\Spf(\overline A)$; then $H$ is a formal subvariety of $G$. Since $J_B\subset I_A$, the augmentation $\varepsilon_A$ induces a continuous morphism of $k$-algebras $\overline\varepsilon:\overline A\to k$.

If (i) holds, then $\Delta_A(I_B)\subset I_B\widehat\otimes1+A\widehat\otimes J_B$, and hence $\Delta_A$ induces, on passing to the quotient, a diagonal morphism $\overline\Delta$. Then $\overline\Delta$ and $\overline\varepsilon$ equip $H$ with a structure of formal submonoid of $G$. Since $G$ is infinitesimal, the same holds for $H$; thus, by Proposition 2.7, $H$ is a formal subgroup of $G$. It then follows from the definition of $G/H$ (cf. 2.4) that (ii) implies the last assertion of the proposition.

On the other hand, it is clear that (ii) implies (i). The proof of the converse occupies paragraphs 5.1.1 to 5.1.5.

\numpara{5.1.1}First consider the following category $\mathscr C$: an object of $\mathscr C$ is a pair $(A,J)$ consisting of a profinite $k$-algebra $A$ and a closed ideal $J$ of $A$; a morphism $\psi:(A,J)\to(A',J')$ of $\mathscr C$ is a continuous homomorphism of $k$-algebras $A\to A'$ which maps $J$ into $J'$. If one associates to $(A,J)$ the pair $(\Spf(A/J),\Spf(A))$, one clearly obtains an antiequivalence from $\mathscr C$ onto the category of pairs $(Z,Y)$ consisting of a formal $k$-variety $Y$ and a formal subvariety $Z$, a morphism $\phi:(Z,Y)\to(Z',Y')$ being a morphism of formal $k$-varieties $Y\to Y'$ which maps $Z$ into $Z'$.

A \emph{cogroup} structure on an object $(A,J)$ of $\mathscr C$ consists of the data of a formal group structure on $\Spf(A)$ such that the following conditions hold (notation of 2.1):
\begin{enumerate}
\item[(1)] $\Delta_A(J)\subset J\widehat\otimes_k A+A\widehat\otimes_k J$;
\item[(2)] $\varepsilon_A(J)=0$;
\item[(3)] $c_A(J)\subset J$.
\end{enumerate}
These conditions also mean that $H=\Spf(A/J)$ is a formal subgroup of $G=\Spf(A)$.\nde{Note that if $(A',J')$ is a second cogroup of $\mathscr C$, corresponding to a pair $H'\subset G'$ of formal $k$-groups, then giving a morphism of cogroups $(A',J')\to(A,J)$ is equivalent to giving a morphism of formal $k$-groups $G\to G'$ which maps $H$ into $H'$.}

Suppose, moreover, that $A$ is \emph{local}, \ie that $\Spf(A)$ is an infinitesimal formal group. Then, if $J\ne A$, conditions (2) and (3) are consequences of (1). Indeed, if $J$ is a closed ideal distinct from $A$, it is contained in the augmentation ideal $I_A$, so (2) holds, and $M=\Spf(A/J)$ is a formal submonoid of $G$. Since $G$ is infinitesimal, it follows from 2.7 that $M$ is a formal subgroup of $G$, \ie condition (3) holds.
% typo?: the asserted implications for general pseudocompact k are retained without new flatness hypotheses.

\numpara{5.1.2}Denote by $\mathbf{Alpg}/k$ the category of ``\emph{graded profinite}'' $k$-algebras: an object of this category consists of the data of a sequence $A_0,A_1,\ldots,A_n,\ldots$ of pseudocompact $k$-modules and a profinite algebra structure on the product $\prod_{n\ge0}A_n$ such that one has $A_n\cdot A_m\subset A_{m+n}$ ($A_n$ is identified with a subset of $\prod_{i\ge0}A_i$ by means of the canonical injection); a morphism $\psi:(A_n)\to(B_n)$ is a sequence of continuous linear maps $\psi_n:A_n\to B_n$ such that one has $\psi_{m+n}(a\cdot a')=\psi_m(a)\cdot\psi_n(a')$ if $a\in A_m$ and $a'\in A_n$.
\begin{definitions}{}
It is clear that two graded profinite $k$-algebras $(A_n)$ and $(B_n)$ have a \emph{coproduct}\nde{We have replaced ``direct sum'' by ``coproduct''.} in $\mathbf{Alpg}/k$, whose $n$th component is the product $\prod_{i=0}^n A_i\widehat\otimes_k B_{n-i}$ of the pseudocompact $k$-modules $A_i\widehat\otimes_k B_{n-i}$. This coproduct will be denoted by $(A_n)\widehat\otimes_k(B_n)$.

Then a cogroup structure on an object $(A_n)$ of $\mathbf{Alpg}/k$ is the data of continuous $k$-linear maps $\Delta_n:A_n\to\prod_{i=0}^n A_i\widehat\otimes_k A_{n-i}$ and $\varepsilon:A_0\to k$ which induce on $\prod_{n\ge0}A_n$ (putting $\varepsilon(A_i)=0$ for $i\ge1$) a cogroup structure in $\mathbf{Alp}/k$.

Finally, for every object $(A,J)$ of $\mathscr C$, one denotes by $\operatorname{Gr}_J(A)$ the graded profinite algebra associated with the filtration of $A$ by the closures $\overline{J^n}$ of the powers of $J$; one therefore has $\operatorname{Gr}_J(A)_n=\overline{J^n}/\overline{J^{n+1}}$ and the multiplication of $\operatorname{Gr}_J(A)$ is induced by that of $A$.
\end{definitions}
\begin{lemma}{}\nde{In the original, the lemma is stated when $k$ is a field, the proof being left to the reader in this case.}
Let $U,V$ be two pseudocompact $k$-modules, with $U$ topologically flat, and let $U=U_0\supset U_1\supset\cdots$ and $V=V_0\supset V_1\supset\cdots$ be two decreasing sequences of closed sub-$k$-modules. Filter the completed tensor product $W=U\widehat\otimes_k V$ by means of the closed submodules
\[
W_n=U_n\widehat\otimes_k V_0+U_{n-1}\widehat\otimes_k V_1+\cdots+U_0\widehat\otimes_k V_n.
\]
Suppose that each $U_i/U_{i+1}$ is topologically flat over $k$ (so that $U/U_n$, hence also $U_n$, is as well, for every $n$). Then, for every $n$, one has an isomorphism
\[
W_n/W_{n+1}\simeq\bigoplus_{i+j=n}(U_i/U_{i+1})\widehat\otimes_k(V_j/V_{j+1}).
\]
\end{lemma}
\begin{proof}
Put $W_{i,j}=U_i\widehat\otimes V_j$ and $\overline W_{i,j}=(U_i/U_{i+1})\widehat\otimes_k(V_j/V_{j+1})$, for every $i,j\ge0$. Let us show by induction on $n$ that the natural map
\[
\pi_n:W_n\longrightarrow\bigoplus_{i+j=n}\overline W_{i,j}
\]
is surjective and that the inclusion $W_{n+1}\subset\Ker(\pi_n)$ is an equality. For $n=0$, the projection
\[
\pi_0:U_0\widehat\otimes_k V_0\longrightarrow(U_0/U_1)\widehat\otimes_k(V_0/V_1)
\]
is surjective and, since $U_0$, $U_0/U_1$ and hence $U_1$ are topologically flat over $k$, one sees that $\Ker(\pi_0)=U_0\widehat\otimes_k V_1+U_1\widehat\otimes_k V_0$ and that, moreover, $U_0\widehat\otimes_k V_1\cap U_1\widehat\otimes_k V_0=U_1\widehat\otimes_k V_1$.

Suppose therefore that $n>0$ and the result is established for $n-1$. Put $M_0=U_0\widehat\otimes_k V_n$ and $S_0=\sum_{i=1}^n U_i\widehat\otimes_k V_{n-i}$. One has $S_0\subset U_1\widehat\otimes_k V_0$ and hence, by the preceding passage applied to $V_0\supset V_n$ instead of $V_0\supset V_1$, one has
\[
M_0\cap S_0\subset U_0\widehat\otimes_k V_n\cap U_1\widehat\otimes_k V=U_1\widehat\otimes_k V_n,
\]
whence one deduces that $M_0\cap S_0=U_1\widehat\otimes_k V_n$. Since $W_n=M_0+S_0$, one obtains a commutative diagram with exact rows, where one has put $U_i'=U_{i+1}$ and $\overline W'_{i,n-1-i}=\overline W_{i+1,n-1-i}$ for $i=0,\ldots,n-1$:
\[
\xymatrix@C=12pt@R=25pt{
0\ar[r]&S_0\ar[r]\ar[d]_{\pi'_{n-1}}&W_n\ar[r]\ar[d]_{\pi_n}&(U_0/U_1)\widehat\otimes_k V_n\ar[r]\ar[d]_p&0\\
0\ar[r]&\bigoplus_{i=0}^{n-1}\overline W'_{i,n-1-i}\ar[r]&\bigoplus_{i=0}^n\overline W_{i,n-1}\ar[r]&\overline W_{0,n}\ar[r]&0.}
\]
% typo?: middle bottom indices i,n-1 are retained as printed.
Then $p$ is surjective, with kernel $(U_0/U_1)\widehat\otimes_k V_{n+1}$. Moreover, by the induction hypothesis applied to the sequence $(U_i')$, $\pi'_{n-1}$ is surjective, with kernel equal to $W_n'=\sum_{i=1}^n W_{i,n+1-i}$. It follows that $\pi_n$ is surjective, and that the inclusion $W_{n+1}\subset\Ker(\pi_n)$ is an equality. This proves the lemma.
\end{proof}
Return to an object $(A,J)$ of $\mathscr C$ and note that, by 0.2.G, the hypothesis that each $\overline{J^n}/\overline{J^{n+1}}$ be topologically flat over $k$ is equivalent to saying that $\operatorname{Gr}_J(A)$ is topologically flat over $k$.
\begin{corollary}{}
Let $\mathscr P$ be the full subcategory of $\mathscr C$ consisting of the objects $(A,J)$ such that $A$ and $\operatorname{Gr}_J(A)$ are topologically flat over $k$. Then the functor $\mathscr P\to\mathbf{Alpg}/k$, $(A,J)\mapsto\operatorname{Gr}_J(A)$ commutes with finite coproducts, hence transforms a cogroup of $\mathscr P$ into a cogroup of $\mathbf{Alpg}/k$.

In particular, if $k$ is a field then, for every cogroup $(A,J)$ of $\mathscr C$, $\operatorname{Gr}_J(A)$ is a cogroup of $\mathbf{Alpg}/k$.\nde{To a pair $H\subset G$ of formal $k$-groups, one therefore associates the ``formal completion $\widehat G_H$ of $G$ along $H$'', which is a formal $k$-group; moreover, we shall see in 5.1.3--5.1.4 that the inclusion $\sigma:H\hookrightarrow\widehat G_H$ possesses a retraction $\pi:\widehat G_H\to H$ and that the formal $k$-group $N=\Ker(\pi)$ is identified, as a formal variety, with the completion of the homogeneous space $G/H$ along the identity section. This will be useful in 5.2.2.}
\end{corollary}
% typo?: the associated graded is described as a formal completion, retained.

\numpara{5.1.3}Identify every profinite $k$-algebra $\Gamma$ with the graded profinite $k$-algebra $(\Gamma_n)_{n\ge0}$ such that $\Gamma_0=\Gamma$ and $\Gamma_n=0$ if $n>0$. In particular, if $(A_n)_{n\ge0}$ is a graded profinite $k$-algebra, we shall consider $A_0$ indifferently as a profinite $k$-algebra or as a graded profinite $k$-algebra. We shall then denote by $\rho:(A_n)\to A_0$ the morphism of $\mathbf{Alpg}/k$ such that $\rho_0=\id_{A_0}$ and $\rho_n=0$ if $n>0$. Similarly, $\tau:A_0\to(A_n)$ will denote the section of $\rho$ such that $\tau_0=\id_{A_0}$ and $\tau_n=0$ if $n>0$.

Every cogroup structure on $(A_n)_{n\in\NN}$ induces a cogroup structure on $A_0$ such that $\rho$ and $\tau$ are homomorphisms of cogroups. In this case, denote by $I_0$ the augmentation ideal of $A_0$ and put $A_n'=A_n/\overline{I_0A_n}$ for every $n\ge0$ (so that $A_0'=k$).

Then $(A_n')_{n\in\NN}$ is a cogroup in $\mathbf{Alpg}/k$ (note that, since $A_0=I_0\oplus k\cdot1$, $I_0\widehat\otimes_k A_n\simeq\overline{I_0A_n}$ is a direct summand of $A_n$, for every $n$). Since $\tau$ is a section of $\rho$, by 2.4.A the cogroup $(A_n)_{n\in\NN}$ is the ``semidirect coproduct'' of $A_0$ and the cogroup $(A_n')_{n\in\NN}$. More precisely, $(A_n')_{n\in\NN}$ is isomorphic, as an object of $\mathbf{Alpg}/k$, to the kernel of the pair:
\[
\xymatrix@C=50pt{(A_n)\ar@<.6ex>[r]^{\tau_1}\ar@<-.6ex>[r]_{(\id\widehat\otimes\rho)\Delta}&(A_n)\widehat\otimes_k A_0}
\]
(where $\Delta:(A_n)\to(A_n)\widehat\otimes(A_n)$ is the comultiplication of $(A_n)$ and $\tau_1(x)=x\widehat\otimes1$), and, identifying $A_n'$ with its image in $A_n$, the map
\[
(A_n'\widehat\otimes_k A_0)\longrightarrow(A_n),\qquad a_n'\widehat\otimes a_0\longmapsto a_n'a_0
\]
is an isomorphism of $\mathbf{Alpg}/k$. (N.B. This is not an isomorphism of cogroups, but $\Delta(A')\subset A\widehat\otimes A'$ and $(\gamma\rho'\widehat\otimes\id)\circ\Delta|_{A'}=\Delta'$, where $\Delta'$ is the comultiplication of $A'$ and $\gamma\rho'$ the projection $A\to A'$, cf. 2.4.A.)

\numpara{5.1.4}Let $(A,J)$ be an object of $\mathscr C$ and $(A_n)=\operatorname{Gr}_J(A)$ the associated object of $\mathbf{Alpg}/k$, \ie $A_n=\overline{J^n}/\overline{J^{n+1}}$ for every $n\ge0$. It is clear that the algebra $\mathscr A=\prod_{n\ge0}A_n$ is generated by $A_0$ and $A_1$, that is, for $n\ge1$, the map
\[
A_1\widehat\otimes_{A_0}A_1\widehat\otimes_{A_0}\cdots\widehat\otimes_{A_0}A_1\longrightarrow A_n\tag{1}
\]
defined by multiplication is surjective.

Suppose, moreover, that $(A,J)$ is a cogroup of $\mathscr C$ and that $A$ and the quotients $\overline{J^n}/\overline{J^{n+1}}$ are flat over $k$. Then, by Corollary 5.1.2, $\operatorname{Gr}_J(A)$ is a cogroup of $\mathbf{Alpg}/k$. Thus, by 5.1.3, if one puts
\[
\begin{split}
A_n'=\biggl\{x\in\overline{J^n}/\overline{J^{n+1}}\ \biggm|\
\Delta(x)-x\widehat\otimes1\in{}&\\[-12pt]
\bigoplus_{i=1}^n(\overline{J^{n-i}}/\overline{J^{n-i+1}})\widehat\otimes(\overline{J^i}/\overline{J^{i+1}})&\biggr\},
\end{split}\tag{2}
\]
then the map $(A_n'\widehat\otimes A_0)\to(A_n)$, $a_n'\widehat\otimes a_n'\mapsto a_n'a_0$ is an isomorphism of $\mathbf{Alpg}/k$.\nde{Let $G,H$ and $\widehat G_H$ be the formal $k$-groups corresponding to $A$, $A_0=A/J$ and $\operatorname{Gr}_J(A)$; then $\tau:A_0\hookrightarrow\operatorname{Gr}_J(A)$ corresponds to a retraction $\pi:\widehat G_H\to H$ of the inclusion $H\hookrightarrow\widehat G_H$, and the preceding passage means that $\widehat G_H$ is the semidirect product of $N=\Ker(\pi)$ by $H$.}
% typo?: a_n' tensor a_n' instead of a_n' tensor a_0 is retained as printed.

Denoting by $I_0$ the augmentation ideal of $A_0$, one deduces from (1) and from the commutative diagram below, where $A_1'^{\widehat\otimes n}$ denotes $A_1'\widehat\otimes_k\cdots\widehat\otimes_k A_1'$ ($n$ factors):
\[
\xymatrix@C=10pt@R=30pt{
A_1\widehat\otimes_{A_0}\cdots\widehat\otimes_{A_0}A_1\ar[d]_m&A_1'^{\widehat\otimes n}\widehat\otimes_k A_0\ar[l]_\sim\ar[d]^{m'\widehat\otimes\id}&(A_1'^{\widehat\otimes n}\widehat\otimes_k I_0)\oplus A_1'^{\widehat\otimes n}\ar[l]_\sim\ar[d]^{m'\widehat\otimes\id\,\oplus m'}\\
A_n&A_n'\widehat\otimes_k A_0\ar[l]_\sim&(A_n'\widehat\otimes_k I_0)\oplus A_n'\ar[l]_\sim}
\]
that the map
\[
m':A_1'\widehat\otimes_k\cdots\widehat\otimes_k A_1'\longrightarrow A_n'\tag{3}
\]
induced by multiplication is surjective; in other words, the profinite $k$-algebra $\mathscr A'=\prod_{n\ge0}A_n'$ is generated by its degree $1$ terms.

\nde{We have detailed the original in what follows, and placed at the end the ``supplement'' $\overline{I_B^n}=B\cap\overline{J_B^n}$ (which is not necessary to establish Proposition 5.1).}Return now to hypothesis (i) of Proposition 5.1: let $G$ be a topologically flat infinitesimal formal $k$-group, $A$ its affine algebra, $I_A$ the augmentation ideal of $A$, $B$ a closed subalgebra of $A$, and $J=\overline{AI_B}$, where $I_B=B\cap I_A$. Denote by $H$ the formal subgroup $\Spf(A/J)$ and by $\pi$ the projection $A\to A/J$. Suppose that $A/\overline{J^n}$ is topologically flat over $k$, for every $n\ge1$, and that $B$ is contained in the kernel $\widetilde B$ of the pair:
\[
\xymatrix@C=55pt{A\ar@<.6ex>[r]^{\tau_1}\ar@<-.6ex>[r]_{(\id_A\widehat\otimes_k\pi)\Delta_A}&A\widehat\otimes_k(A/J).}
\]
Let $(A_n)=\operatorname{Gr}_J(A)$, let $I_0=I_A/J$ be the augmentation ideal of $A_0=A/J$, and define $(A_n')_{n\in\NN}$ as in (2) above. Denote by $(B_n)_{n\in\NN}$ (resp. $(\widetilde B_n)_{n\in\NN}$) the object of $\mathbf{Alpg}/k$ associated with the filtration of $B$ (resp. $\widetilde B$) induced by that of $A$, \ie defined by the ideals $B\cap\overline{J^n}$ (resp. $\widetilde B\cap\overline{J^n}$). Then it is clear that $B_n\subset\widetilde B_n\subset A_n'$ for every $n$, and that
\[
\mathscr B=\prod_{n\ge0}B_n\subset\widetilde{\mathscr B}=\prod_{n\ge0}\widetilde B_n
\]
are subalgebras of $\mathscr A'=\prod_{n\ge0}A_n'$.

On the other hand, $J$ (resp. $\overline{J^2}$) is the image in $A$ of $I_B\widehat\otimes_k A$ (resp. $\overline{I_B^2}\widehat\otimes_k A$). Consequently, the map
\[
(I_B/\overline{I_B^2})\widehat\otimes_k A\longrightarrow J/\overline{J^2}=A_1\simeq A_1'\widehat\otimes_k A_0
\]
is surjective, and factors through $(I_B/\overline{I_B^2})\widehat\otimes_k A_0$. Since $A_0=k\cdot1\oplus I_0$, one deduces from the commutative exact diagram:
\[
\xymatrix@C=22pt@R=22pt{
(I_B/\overline{I_B^2})\widehat\otimes_k A_0\ar[d]&(I_B/\overline{I_B^2})\widehat\otimes_k I_0\oplus(I_B/\overline{I_B^2})\ar[l]_\sim\ar[d]\\
J/\overline{J^2}\ar[d]&A_1'\widehat\otimes_k I_0\oplus A_1'\ar[l]_\sim\ar[d]\\
0&0}
\]
that $A_1'$ is the image of $I_B/\overline{I_B^2}$, so that one has $B_1=A_1'$. Since $\mathscr A'$ is generated by $A_1'$, it follows that, for every $n\ge1$, one has
\[
A_n'\subset B_n\subset\widetilde B_n\subset A_n',
\]
whence $B_n=\widetilde B_n=A_n'$.\nde{With the notation of N.D.E. (149), this implies that the formal completion of $G/H$ along the identity section (whose affine algebra is $\prod_n B_n$) is isomorphic, as a formal variety, to the formal $k$-group $N$.}

Finally, since the formal group $G$ is \emph{infinitesimal}, one has $\mathfrak r(A)=I_A$ and hence the ideals $\overline{I_A^n}$ tend to $0$ (cf. 0.1.2); a fortiori, the ideals $\overline{J^n}$ tend to $0$, and hence the induced filtrations on $\widetilde B$ and $B$ are separated. Moreover, since $B$ is a \emph{closed} subalgebra of $A$, it is complete for the topology defined by the ideals $B\cap\overline{J^n}$. Consequently, it follows from \cite{CA}, § V.7, Lemma 1 (see also 5.1.5 below) that $B=\widetilde B$. This completes the proof of Proposition 5.1.

One has, moreover, the following supplement. For every $n$, $\overline{J^n}=\overline{AI_B^n}$ is the image in $A$ of $A\widehat\otimes_B\overline{I_B^n}$ and also of $A\widehat\otimes_B(B\cap\overline{J^n})$. Now, by hypothesis, the affine algebra $A/J$ of the formal subgroup $H$ is topologically flat over $k$. Thus, by Theorem 2.4, the morphism $G=\Spf(A)\to G/H=\Spf(B)$ is \emph{surjective and topologically flat}; one therefore has
\[
A\widehat\otimes_B\overline{I_B^n}=\overline{J^n}=A\widehat\otimes_B(B\cap\overline{J^n}),
\]
and this implies that $\overline{I_B^n}=B\cap\overline{J^n}$ for every $n$. This also follows from the fact that the maps
\[
\overline{I_B^n}/\overline{I_B^{n+1}}\longrightarrow A_i'=(B\cap\overline{J^n})/(B\cap\overline{J^{n+1}})
\]
are surjective, and from 5.1.5 (ii) below, applied to $B_n'=\overline{I_B^n}$ and $B_n=B\cap\overline{J^n}$.
% typo?: A_i' instead of A_n' is retained as printed.

\begin{lemma}{5.1.5}\label{VIIB.5.1.5}\nde{The original stated only point (ii); for the reader's convenience, we have stated Lemma 1 of \cite{CA}, § V.7, in (i).}
(i) Let $M$ and $N$ be two abelian groups filtered by decreasing sequences of subgroups $(M_n)_{n\in\ZZ}$ and $(N_n)_{n\in\ZZ}$. Suppose that the union of the $M_n$ (resp. $N_n$) equals $M$ (resp. $N$), that the intersection of the $M_n$ (resp. $N_n$) is zero, and that $M$ is complete for the topology defined by the $M_n$. Let $f:M\to N$ be a morphism of filtered groups.
\begin{enumerate}
\item[a)] If $f$ induces a surjection of the associated gradeds, then $f$ is a surjection and $N$ is complete for the topology defined by the $N_n$.
\item[b)] If $f$ induces an injection of the associated gradeds, then $f$ is an injection.
\end{enumerate}
(ii) Let $B$ be an abelian group, and let $B=B_0'\supset B_1'\supset\cdots$ and $B=B_0\supset B_1\supset\cdots$ be two separated filtrations of $B$ by subgroups such that $B_n'\subset B_n$ for every $n$. Suppose that $B$ is complete for the topology defined by the filtration $(B_n')$.

If the map $B_i'/B_{i+1}'\to B_i/B_{i+1}$ is surjective for every $i$, then $B_n'=B_n$ for every $n$.
\end{lemma}
Indeed, (i) is Lemma 1 of \cite{CA}, § V.7 (see also \cite{BAC}, III, § 2.8), and (ii) follows by taking $M=B_n'\supset B_{n+1}'\supset\cdots$ and $N=B_n\supset B_{n+1}\supset\cdots$.

\numpara{5.2}\emph{Throughout the remainder of Section 5, $k$ denotes a perfect field of characteristic $p>0$}.

We put $\overline\NN=\NN\cup\{\infty\}$. If $B$ is a profinite $k$-algebra and $r\in\overline\NN$, we denote by $((x^{p^r}))_{x\in\mathfrak r(B)}$ the closed ideal of $B$ generated by the elements $x^{p^r}$, where $x$ ranges over the radical $\mathfrak r(B)$ of $B$. If $r=\infty$, we use the same notation, agreeing that $((x^{p^\infty}))_{x\in\mathfrak r(B)}$ is the zero ideal. In both cases, $B_r$ denotes the quotient $B/((x^{p^r}))_{x\in\mathfrak r(B)}$.

We say that $B$ has \emph{height $\le r$} if $((x^{p^r}))_{x\in\mathfrak r(B)}$ is the zero ideal; if this holds and $r$ is finite, we say that $B$ has \emph{finite height}.

Consider in particular the case where $B$ is of the form $k[[\omega]]$, $\omega$ being a pseudocompact $k$-vector space (cf. 1.2.5).\nde{In the original, the author uses ``linearly compact vector space'', which is equivalent to ``pseudocompact vector space'' (cf. \cite{BAC}, § III.2, Exercises 15 a), 19 a) and 20 d)). We have preferred to retain the terminology ``pseudocompact'', used up to this point.} We then say that $B$ is a \emph{formal power series algebra} and that $B_r$ is a \emph{truncated formal power series algebra} ($r\in\overline\NN$; we therefore agree to say that $B=B_\infty$ is also ``truncated''). If $B=k[[\omega]]$, we also write $((x^{p^r}))_{x\in\omega}$ instead of $((x^{p^r}))_{x\in\mathfrak r(B)}$.
\begin{notation}{}
Let $\omega$ be a pseudocompact $k$-vector space filtered by an increasing sequence of closed vector subspaces
\[
0=\omega_0\subset\omega_1\subset\omega_2\subset\omega_3\subset\cdots.
\]
(a) The closed ideal of $k[[\omega]]$ generated by the elements $x^{p^r}$, where $r$ ranges over $\NN$ and $x$ ranges over $\omega_r$, will be denoted by $((x^{p^r}))_{r\in\omega_r}$.
% typo?: r in omega_r, rather than x in omega_r, is printed in this notation, retained.

(b) On the other hand, we shall denote by ${}_r\omega$ the filtered pseudocompact vector space such that
\[
{}_r\omega_i=\omega_i\quad\text{if }i<r,\qquad{}_r\omega_i=\omega\quad\text{if }i\ge r.
\]
\end{notation}
\begin{theorem}{}\label{VIIB.5.2}\textup{(Dieudonné--Cartier).}
Let $H\to G$ be a monomorphism of infinitesimal formal groups over a perfect field $k$ of characteristic $p>0$. Let $B$ be the affine algebra of the homogeneous space $G/H$ and suppose that one of the following three conditions holds:\nde{On the one hand, we have replaced $H\backslash G$ by $G/H$, and likewise in the proof; on the other hand, we have added condition (iii).}
\begin{enumeratei}
\item $B$ has finite height (this holds in particular if $G$ has finite height).
\item $B$ is a complete noetherian local ring.
\item $B$ is a reduced ring.
\end{enumeratei}
Then $B$ is isomorphic to the completed tensor product of a finite family of truncated formal power series algebras.
\end{theorem}
The proof of this theorem occupies paragraphs 5.2.1 to 5.2.5.

\numpara{5.2.1}Let $A$ be the affine algebra of $G$, $I_A$ its augmentation ideal, and $I=B\cap I_A$. By 2.4, one has $H=\Spf(A/\overline{AI})$ and $B=\{x\in A\mid\Delta(x)-1\widehat\otimes x\in A\widehat\otimes\overline{AI}\}$. Put $\omega=I/\overline{I^2}$. One denotes by $I_r$ the closed subideal of $I$ consisting of the $x$ such that $x^{p^r}=0$, and by $\omega_r$ the canonical image of $I_r$ in $\omega$. We shall prove:
% typo?: 1 tensor x, rather than x tensor 1, is printed here, retained.
\begin{proposition}{}
If there exists a continuous section $\sigma:\omega\to I$ of the projection $I\to I/\overline{I^2}$ such that $\sigma(\omega_r)\subset I_r$ for every $r$, then $B$ is isomorphic to $k[[\omega]]/((x^{p^r}))_{x\in\omega_r}$.
\end{proposition}
Such a section indeed extends to a continuous morphism $k[[\omega]]\to B$, which factors through $B'=k[[\omega]]/((x^{p^r}))_{x\in\omega_r}$. We prove in 5.2.2 to 5.2.5 that the morphism
\[
\phi:B'=k[[\omega]]/((x^{p^r}))_{x\in\omega_r}\longrightarrow B
\]
thus obtained is an isomorphism.

\numpara{5.2.1.A}\label{VIIB.5.2.1.A}\nde{We have added paragraphs 5.2.1.A and 5.2.1.B.}For each $r\in\NN$, put $\mathcal I_r=\overline{I^2}+I_r$, so that $\mathcal I_r/\overline{I^2}\simeq\omega_r$; one then has a commutative diagram with exact rows
\[
\xymatrix@C=20pt@R=24pt{
0\ar[r]&\mathcal I_{r-1}\ar[r]\ar[d]&\mathcal I_r\ar[r]\ar[d]&\mathcal I_r/\mathcal I_{r-1}\ar[r]\ar[d]^\wr&0\\
0\ar[r]&\omega_{r-1}\ar[r]&\omega_r\ar[r]&\omega_r/\omega_{r-1}\ar[r]&0}
\]
and since $k$ is a field, the rows are split: one can extend a pseudobasis $\mathscr B_{r-1}$ of $\mathcal I_{r-1}$ to a pseudobasis $\mathscr B_{r-1}\cup\mathscr B_r'$ of $\mathcal I_r$, and then the closed subspace $S_r$ with pseudobasis $\mathscr B_r'$ is a complement to $\mathcal I_{r-1}$ in $\mathcal I_r$, and the projection $\pi:I\to\omega$ induces an isomorphism from $S_r$ onto a complement $\omega_r'$ to $\omega_{r-1}$ in $\omega_r$. Denote by $\mathcal I_\infty$ the closed ideal $\overline{\bigcup_r\mathcal I_r}$; it likewise admits a complement $S_\infty$ in $I$, and $\pi$ induces an isomorphism from $\mathcal I_\infty/\overline{I^2}$ (resp. from $S_\infty$) onto the closure $\omega_\infty$ of the union of the $\omega_r$ (resp. onto a complement $\omega_\infty'$ to $\omega_\infty$ in $\omega$). Denote by $\eta$ the isomorphism $S_\infty\isomto\omega_\infty'$. One then obtains continuous linear maps:
\[
\xymatrix@C=25pt@R=25pt{
\overline{I^2}\times S_\infty\times\widehat{\bigoplus}_r S_r\ar[r]^\phi\ar[d]_{\eta\times\theta}&I\\
\omega=\omega_\infty'\times\omega_\infty&}
\]
where $\widehat{\bigoplus}_r S_r$ is the direct sum of the $S_r$ in $\mathbf{PC}(k)$, \ie $(\prod_r S_r^\dagger)^*$ (cf. N.D.E. (16) of 0.2.2), and where $\theta:\widehat{\bigoplus}_r S_r\to\omega_\infty$ is induced by the maps $S_r\isomto\omega_r'\hookrightarrow\omega_\infty$.
